% !TEX root = paper-node2vec.tex

% Aditya's version
% \hide{
% Many important tasks in network analysis involve some kind of prediction over nodes and edges. In a typical node classification task, we are interested in predicting the most probable class(es) for nodes in the network. For example, in a social network such as Facebook we might be interested in predicting topics of interest for an individual user to suggest relevant page recommendations. Similary, in link prediction, we wish to predict whether any two given nodes in a network have an edge connecting them. This is useful in a wide variety of domains, for instance, in genomics, it helps us discover novel interactions between genes; in academic collaboration networks, it can identify potential researchers with similar interests etc.

% At the most fundamental level, any supervised learning algorithm over networks requires a feature vector representation for the nodes and edges. A naive solution is to hand-engineer domain specific features based on expert knowledge and augment these with some standard network features such as centrality measures for nodes, Adamic-Adar score~\cite{adamicadar} for edges etc. Even if one discounts the explicit reliance on experts for both domain and network features, limiting the learning algorithms to consider a pre-defined set of network features derived from earlier research is tedious and restrictive. An alternative approach is to turn the problem on its head: can we learn feature representations themselves? In addition to taking us closer to truly automated machine learning, learned feature representations are effective in expressing latent explanatory variables that generalize across prediction tasks. The challenge in feature learning is defining an objective function which involves a trade-off in balancing computational efficiency and statistical accuracy. On one side of the spectrum of algorithms, one could directly optimize the same objective as we have for a downstream prediction task. This is typical in deep networks with many hidden layers, and results in good accuracy, but at the cost of extremely high training time requirements due to a blowup in the number of parameters to be learned. At the other extreme, the objective function is defined to be independent of the downstream task in a purely unsupervised learning setting. This makes the optimization computationally tractable and with a carefully designed objective, can outperform the former approach in balancing the accuracy-speed trade-off. 

% However, current techniques fail to satisfactorially fulfill the criteria required for unsupervised feature learning to work in networks. Classic approaches based on linear and non-linear dimensionality reduction techniques such as Principal Component Analysis (PCA), Multi-Dimensional Scaling (MDS) and their novel extensions~\cite{tenenbaum-science2000, roweis-science2000, belkin-nips2001} invariably involve eigendecomposition of a representative data matrix which is computationally intractable for large real-world networks. In contrast, the widespread success of neural networks sheds light on a radically different approach for unsupervised feature learning. Empirically, it has been observed that a neural network with only a single hidden layer not only guarantees tractability, but at the same time is able to learn rich representations through non-linear transformations over a large number of node-neighbourhood samples from the underlying network. Recent attempts in this direction specifically for network analysis~\cite{deepwalk, line} propose algorithms that overcome the computational bottleneck, but are largely insensitive to patterns unique to networks. For instance, nodes in networks could be classified based on communities they belong to (homophily) whereas in some cases, the classification could be based on the structural roles nodes play in the network (structural equivalence). From a networks perspective, the absence of a rigorous analysis of the design choices involved in these algorithms and consequently, their effect on the learned representations is a major limiting drawback in generalizing these methods across domains.

% In this paper, we propose \nodevec, a semi-supervised algorithm for scalable feature learning in networks based on single layer neural networks. Our objective function is a custom likelihood estimate in line with prior work in natural language processing for unsupervised feature learning~\cite{word2vec} generalized to graphs and intuitively, the optimization is over mappings that project nodes to a $d$-dimensional feature space preserving neighbourhoods. We use a Monte Carlo inference approach for learning, and our key contribution is defining neighbourhoods in terms of cooccuring node samples obtained effciently by simulating biased random walks on the underlying network. Specifically, through the \textit{network} bias, our walks express a preference for nodes which exhibit a high relative importance, measured by an eigenvector centrality and normalized for degree. Secondly, we also incorporate an explicit \textit{search} bias term which correlates classic search procedures with specific equivalences (homophilic and/or structural) observed in networks. Our algorithm is flexible, giving us control over the search space we explore in contrast to rigid search procedures in prior work~\cite{deepwalk, line}, and consequently, can discover the full spectrum of equivalences observed in networks. The parameters governing our search strategy have an intuitive interpretation, and hence, could be set directly based on the underlying domain. Our experiments show that they can also be inferred cheaply from a tiny fraction of labelled data and generalize well at test time.  

% Lastly, for the first time in our knowledge, we show how feature representations of individual nodes learned using neural networks can be extended to pairs of nodes (including edges). We use a bootstrapping approach to generate feature representations for pairs of nodes by applying simple binary operations over the individual nodes constituting the pair. This lends \nodevec to prediction tasks involving edges as well. Our experiments focus on the two most popular prediction tasks in networks. In the first set of experiments, we do multi-label classification where every node is assigned one or more class labels from a finite set of classes. Our second set of experiments is based on link prediction where we predict the existence of an edge given any two pairs of nodes. In all our experiments, we contrast \nodevec against recent state-of-the-art feature learning algorithms~\cite{deepwalk, line} on real world networks from diverse domains such as social networks, systems biology etc. We observe a performance gain upto x\% on the multi-label classification and upto y\% on link prediction over the closest benchmark.

% Overall our paper makes the following contributions:
% \begin{enumerate}
% \item We propose and release \nodevec \footnote{Available at TODO}, an efficient scalable algorithm for feature learning in networks based on neural embeddings.
% \item We adapt \nodevec to established principles in network science, particularly the role of eigenvector centralities and provide a mechanism to tune the search procedure in our algorithm to additional equivalences observed in networks.
% \item We extend \nodevec and more generally, unsupervised feature learning techniques based on neural embeddings from nodes to pairs of nodes for edge-based prediction tasks.
% \item We empirically demonstrate the success of \nodevec on tasks incolving multi-label classification over nodes and link prediction on real-world datasets derived from diverse domains.
% \end{enumerate}

% The rest of the paper is structured as follows. In Section 2, we briefly survey related work in feature learning in networks. We present the technical details for feature learning using \nodevec in Section 3, including methods for generalization of node embeddings to features for pairs of nodes. This is followed by empirical evaluations on node and edge prediction tasks on several real-world datasets in Section 4. Finally, we conclude and discuss potential directions for future work in Section 5.
% } % Hide

%%% JURE

Many important tasks in network analysis involve predictions over nodes and edges. In a typical node classification task, we are interested in predicting the most probable labels of nodes in a network~\cite{tsoumakas2006multi}. For example, in a social network, we might be interested in predicting interests of users, or in a protein-protein interaction network we might be interested in predicting functional labels of proteins~\cite{radivojac2013large,yang-www2011}. Similarly, in link prediction, we wish to predict whether a pair of nodes in a network should have an edge connecting them~\cite{liben2007link}. Link prediction is useful in a wide variety of domains; for instance, in genomics, it helps us discover novel interactions between genes, and in social networks, it can identify real-world friends~\cite{srw,vazquez2003global}.

Any supervised machine learning algorithm requires a set of informative, discriminating, and independent features.
In prediction problems on networks this means that one has to construct a feature vector representation for the nodes and edges. A typical solution involves hand-engineering domain-specific features based on expert knowledge. Even if one discounts the tedious effort required for feature engineering, such features are usually designed for specific tasks and do not generalize across different prediction tasks.


An alternative approach is to 
%turn the problem on its head and 
 {\em learn} feature representations by solving an optimization problem~\cite{bengio-pami2013}. 
%Learned feature representations are effective in producing latent representations that generalize across prediction tasks. 
The challenge in feature learning is defining an objective function, which involves a trade-off in balancing computational efficiency and predictive accuracy. On one side of the spectrum, one could directly aim to find a feature representation that optimizes performance of a downstream prediction task. 
While this supervised procedure results in good accuracy, it comes at the cost of high training time complexity due to a blowup in the number of parameters that need to be estimated. At the other extreme, the objective function can be defined to be independent of the downstream prediction task and the representations can be learned in a purely unsupervised way. This makes the optimization computationally efficient and with a carefully designed objective, it results in task-independent features that closely match task-specific approaches in predictive accuracy~\cite{word2vec, pennington-emnlp2014}.
%former approach in balancing the accuracy-speed trade-off.

However, current techniques fail to satisfactorily define and optimize a reasonable objective required for scalable unsupervised feature learning in networks. Classic approaches based on linear and non-linear dimensionality reduction techniques such as Principal Component Analysis, Multi-Dimensional Scaling and their extensions~\cite{belkin-nips2001, roweis-science2000, tenenbaum-science2000, yan2007graph} optimize an objective that transforms a representative data matrix of the network such that it maximizes the variance of the data representation. Consequently, these approaches invariably involve eigendecomposition of the appropriate data matrix which is expensive for large real-world networks. Moreover, the resulting latent representations give %explaining network behaviour and 
poor performance on various prediction tasks over networks. 

Alternatively, we can design an objective that seeks to preserve local neighborhoods of nodes. The objective can be efficiently optimized using stochastic gradient descent (SGD) akin to backpropogation on just single hidden-layer feedforward neural networks.
 %Empirically, it has been observed that a neural network with only a single hidden layer is not only computationally efficient, but is also able to learn rich representations through non-linear transformations over a large number of node-neighborhood samples from the underlying network~\cite{word2vec}. 
Recent attempts in this direction~\cite{deepwalk, line} propose 
%algorithms that overcome the computational bottleneck, 
efficient algorithms but %are largely insensitive to patterns unique to networks. 
rely on a rigid notion of a network neighborhood, which results in these approaches
being largely insensitive to connectivity patterns unique to networks.
%
Specifically, nodes in networks could be organized based on communities they belong to (\ie, \emph{homophily}); in other cases, the organization could be based on the structural roles of nodes in the network (\ie, \emph{structural equivalence})~\cite{fortunato09community,henderson-kdd2012, yang14ieee}. 
% \sout{These roles could allude to different interpretations; in some cases, the roles are domain-dependent, or they could even be generally defined in terms of network structure such as hubs, bridges, etc.}
For instance, in Figure~\ref{fig:bfsdfs}, we observe nodes $u$ and $s_1$ belonging to the same tightly knit community of nodes, while the nodes $u$ and $s_6$ in the two distinct communities share the same structural role of a hub node. Real-world networks commonly exhibit a mixture of such equivalences. Thus, it is essential to allow for a flexible algorithm that can learn node representations obeying both principles: ability to learn representations that embed nodes from the same network community closely together, as well as to learn representations where nodes 
%with similar structural roles get embedded together
that share similar roles have similar embeddings. This would allow feature learning algorithms to generalize across a wide variety of domains and prediction tasks.

\xhdr{Present work}
We propose \nodevec, a semi-supervised algorithm for scalable feature learning in networks. We optimize
a custom graph-based objective function using SGD motivated by prior work on natural language processing 
%for unsupervised feature learning
~\cite{word2vec}.
Intuitively, our approach returns feature representations that maximize the likelihood of preserving network neighborhoods of nodes in a $d$-dimensional feature space. 
We use a 2$^{\textrm{nd}}$ order random walk approach to generate (sample) network neighborhoods for nodes.

Our key contribution is in defining a flexible notion of a node's network neighborhood. By choosing an appropriate notion of a neighborhood, \nodevec can learn representations that organize nodes based on their network roles and/or communities they belong to. We achieve this by developing a family of biased random walks, which efficiently explore diverse neighborhoods of a given node. 
% Secondly, we also incorporate an explicit \textit{search} bias term, which biases the random walk towards specific node equivalence classes (homophilic and/or structural). 
The resulting algorithm is flexible, giving us control over the search space through tunable parameters, in contrast to rigid search procedures in prior work~\cite{deepwalk, line}. Consequently, our method generalizes prior work and can model the full spectrum of equivalences observed in networks. The parameters governing our search strategy have an intuitive interpretation and bias the walk towards different network exploration strategies. These parameters can also be learned directly using a tiny fraction of labeled data in a semi-supervised fashion.

We also show how feature representations of individual nodes can be extended to pairs of nodes ({\em i.e.}, edges). In order to generate feature representations of edges, we compose the learned feature representations of the individual nodes using simple binary operators. This compositionality lends \nodevec to prediction tasks involving nodes as well as edges. 

Our experiments focus on two common prediction tasks in networks: a multi-label classification task, where every node is assigned one or more class labels and a link prediction task, where we predict the existence of an edge given a pair of nodes.
We contrast the performance of \nodevec with state-of-the-art feature learning algorithms~\cite{deepwalk, line}.
We experiment with several real-world networks from diverse domains, such as social networks, information networks, as well as networks from systems biology. Experiments demonstrate that \nodevec outperforms state-of-the-art methods by up to 26.7\% on multi-label classification and up to 12.6\% on link prediction. The algorithm shows competitive performance with even 10\% labeled data and is also robust to perturbations in the form of noisy or missing edges. Computationally, the major phases of \nodevec are trivially parallelizable, and it can scale to large networks with millions of nodes in a few hours.

% \jure{We want to say more here. Talk about scalability, performance vs. amount of training data, efficiency of sampling, etc.}

Overall our paper makes the following contributions:
\begin{enumerate}[noitemsep,nolistsep]
	\item We propose \nodevec,  
	an efficient scalable algorithm for feature learning in networks that efficiently optimizes a novel network-aware, neighborhood preserving objective using SGD.
	\item We show how \nodevec is in accordance with established principles in network science, providing flexibility in discovering representations conforming to different equivalences.
	\item We extend \nodevec and other feature learning methods based on neighborhood preserving objectives, from nodes to pairs of nodes for edge-based prediction tasks.
	\item We empirically evaluate \nodevec for multi-label classification and link prediction on several real-world datasets.
\end{enumerate} 

The rest of the paper is structured as follows. In Section~\ref{sec:related}, we briefly survey related work in feature learning for networks. We present the technical details for feature learning using \nodevec in Section~\ref{sec:proposed}.
%, including methods for generalization of feature embeddings for nodes to features for pairs of nodes. 
In Section~\ref{sec:experiments}, we empirically evaluate \nodevec on prediction tasks over nodes and edges on various real-world networks 
%along with
and assess the parameter sensitivity, perturbation analysis, and scalability aspects of our algorithm. 
We conclude with a discussion of the \nodevec framework and highlight some promising directions for future work in Section~\ref{sec:discussion}. Datasets and a reference implementation of \nodevec are available on the project page: \url{http://snap.stanford.edu/node2vec}.
